\section{Summary of Findings}

This dissertation set out to determine whether the geographic location a user discloses shapes the mental health guidance produced by contemporary large language models, and to do so using a methodology in which every automated outcome measure is independently validated before being used to support an inference. The answer to the substantive question is yes, and in the direction that disadvantages the users who would most benefit from reliable guidance. The answer to the methodological question is that independent validation is not optional: this study found, and corrected, a measurement failure that would have produced a published finding pointing in the wrong direction.

Three pre-registered hypotheses are supported at the level tested. Response actionability increases monotonically with the income category of the stated location ($+0.614$ on a 0--5 scale, High vs Low, $p < 0.0001$; validated $F_1 = 0.744$). Response localisation increases with the documented availability of a national crisis service ($+0.353$, Established vs None documented, $p < 0.001$; validated $F_1 = 0.750$). The probability of receiving a recommendation to seek religious or spiritual support varies sharply by WHO region, with Africa and the Eastern Mediterranean above every other region and no positive cases observed in Europe (validated $F_1 = 0.889$, 95\% CI [0.678, 0.934]). The H3 model adjusts for model identity and scenario but not income, so this regional association is not claimed to be income-independent.

An unplanned fourth finding, emerging from auditing the raw response data rather than from the pre-registered design, identifies a pattern of visibly fabricated crisis-contact information that concentrates in the locations with the least documented infrastructure: 34.4\% of crisis contacts offered to users in countries with no documented national crisis service display recognisably synthetic digit patterns (fictional-exchange blocks, sequential runs, repeated digits), against 1.1\% in countries with an established national line.

\section{Answers to the Research Questions}

\textbf{RQ1} is answered affirmatively: actionability varies with income category, the gradient is monotone, and the effect survives an analysis conducted on the twenty city-level means only (Spearman $\rho = 0.848$ across city means, $p < 0.0001$).

\textbf{RQ2} is answered affirmatively after correction of a measurement defect that would otherwise have produced the opposite conclusion. Localisation increases with documented crisis-service availability. The finding that the corrected and uncorrected measures point in opposite directions is itself an empirical contribution to the methodology of LLM auditing.

\textbf{RQ3} is answered affirmatively as a regional association: WHO region predicts the probability of religious framing after adjustment for model identity and scenario, with the pattern concentrated in AFR and EMR and absent in EUR. Because income is not included in the fitted H3 model, the analysis does not establish that the association is independent of income. The finding nevertheless documents a systematic location-linked difference made without any user-stated religious preference.

\section{Directions for Future Work}

Three directions are prioritised for future work, ordered by the degree to which they would strengthen the chain of inference from this study.

First, the crisis-service reference data should be completed against a verified, current institutional directory --- the WHO Mental Health Atlas originally proposed --- for all twenty countries. Five countries were directly re-verified during this dissertation; the remaining fifteen retain desk-research classifications. This is the weakest link in both the H2 analysis and the fabrication finding, and it is the highest-value remaining task.

Second, a no-location control condition should be added: the same eight scenarios posed with no city or country stated. This would allow the deviation from a location-free baseline to be estimated directly, distinguishing between ``models know less about lower-income settings'' and ``models withhold information when a lower-income location is stated.'' These are different claims with different policy implications, and the present design cannot separate them.

Third, expanding beyond twenty cities would address the most important bound on generalisability. Twenty cities imposes an effective sample size of twenty for every geographic contrast. A design of sixty or more cities, balanced across income tiers and WHO regions, would allow city-income matching --- comparing cities with similar income but different crisis-service infrastructure, and vice versa --- which is the analysis that separates correlation from mechanism.

\section{Concluding Remarks}

Large language models are already consulted for mental health guidance by users for whom no comparable alternative exists. This study finds that the guidance they provide is not geographically neutral: it is measurably less actionable and less localised in lower-resource settings, culturally framed differently by region without reference to individual user preferences, and in a meaningful proportion of cases in the least-served settings includes fabricated crisis-contact information that would fail a person in acute distress. These findings were established using validated outcome measures, confirmed by three alternative specifications, and documented with full transparency including the one measure that failed validation and would have supported a backwards conclusion.

The methodological contribution may prove as durable as the substantive findings. LLM audit studies that rely on automated content coding without independent human validation face a specific and underappreciated risk: the measure may agree with itself while measuring the wrong thing. This study shows that risk is not hypothetical, and that a validation protocol capable of catching directional failures --- not merely calibrating thresholds --- is both feasible within a dissertation-scale project and necessary for defensible inference.
